% ch06-systems.tex -- pipelining the core, hazards shown honestly
\chapter{Pipelining: One Per Cycle (Usually)}
\label{ch:pipe06}

The single-cycle core of ch05 was honest and slow: one instruction
per very long clock, every instruction priced at \texttt{lw}'s
journey through five hardware stops. This chapter cuts that path
into five stages with four pipeline registers --- IF, ID, EX, MEM,
WB --- and discovers the machine's first real lie. The pipeline
retires one instruction per cycle whether the instructions depend
on each other or not. On independent code it is correct and fast;
on dependent code it is fast and \emph{wrong}: our testbench shows
\texttt{a1} come back 3 instead of 8 because the second
\texttt{addi} reads \texttt{a0} three cycles before it is written.
Hazards are not an edge case --- they are the price of admission.
This chapter builds the broken machine on purpose, measures the
breakage, and hands the fix to ch07.

\section{Why pipeline}
\label{sec:why06}
\index{pipelining}\index{throughput}\index{latency}

Two numbers, not one. \emph{Latency} is how long one instruction
takes end to end --- pipelining does not improve it (pipeline
registers add a little). \emph{Throughput} is how many finish per
unit time --- pipelining multiplies it by the stage count, ideally.
Ch05's $5\tau$ clock becomes a $\tau$ clock; five instructions in
flight finish one per $\tau$ instead of one per $5\tau$. The
\texttt{lw} tax disappears not because \texttt{lw} got cheaper but
because its slack is now occupied by four other instructions.

\begin{figure}[t]
\centering
\begin{tikzpicture}[font=\scriptsize, >=latex]
% laundry timeline: single-cycle vs pipelined
\node[font=\footnotesize\bfseries] at (6,5.6)
  {same 5 instructions, two machines};
% single cycle row
\node[left, font=\tiny] at (0.4,4.9) {single-cycle};
\foreach \i in {0,...,4}{
  \draw[thick, fill=accentsoft] (0.6+\i*1.0,4.5) rectangle (1.55+\i*1.0,5.15);}
\node[font=\tiny] at (3.1,4.25) {each box = one slow clock};
% pipeline lanes
\node[left, font=\tiny] at (0.4,2.6) {pipelined};
\foreach \lane/\lbl in {0/IF,1/ID,2/EX,3/MEM,4/WB}{
  \node[left, font=\tiny] at (6.9,3.4-\lane*0.5) {\lbl};}
\foreach \i in {0,...,4}{
  \foreach \lane in {0,...,4}{
    \pgfmathtruncatemacro{\s}{\i+\lane}
    \draw[thick, fill=accentsoft] (7.0+\s*0.45,3.4-\lane*0.5)
      rectangle (7.45+\s*0.45,3.9-\lane*0.5);}}
\node[font=\tiny] at (8.9,0.4) {instr $n$ in lane $n$ --- five at once};
\node[align=left, text=slate, text width=11cm] at (6.5,-0.7)
  {Left: 5 instructions in 5 long periods. Right: the same five
   overlap --- the pipeline finishes one instruction per short
   period once full. Latency per instruction is unchanged; the
   throughput is what multiplies.};
\end{tikzpicture}

\caption[Laundry timeline]{Same five instructions on both machines.
Single-cycle serialises; the pipeline overlaps --- throughput, not
latency, is what changed.}
\label{fig:timeline06}
\end{figure}

\section{The five stages and what the registers carry}
\label{sec:stages06}
\index{pipeline register}\index{IF/ID/EX/MEM/WB}

\begin{example}[cut the path, keep the semantics]\label{ex:regs06}
\texttt{ch06\_rv32i\_pipe\_top.v} reuses ch05's decoder, regfile and
ALU unmodified --- the stages are the same combinational blocks,
separated by registers. \texttt{IF/ID} carries the instruction word
and \texttt{pc+4}; \texttt{ID/EX} carries the read operands, the
immediate, the destination index, and every control bit decode
produced; \texttt{EX/MEM} carries the ALU result, the store value,
and the surviving control bits; \texttt{MEM/WB} carries whichever
of ALU-result / memory-data / pc+4 the writeback needs. Nothing is
forgotten: a control bit that dies early is a bug with a wire width.
\end{example}

\begin{figure}[t]
\centering
\begin{tikzpicture}[font=\scriptsize, >=latex]
% 5-stage datapath with pipeline registers
\foreach \x/\t/\c in {
  0.8/{IF}/{accentsoft}, 2.9/{ID}/{goodbg}, 5.0/{EX}/{amberbg},
  7.1/{MEM}/{graybg}, 9.2/{WB}/{accentsoft}}{
  \draw[thick, fill=\c] (\x-0.75,2.9) rectangle (\x+0.75,4.1);
  \node[font=\footnotesize] at (\x,3.5) {\t};}
\foreach \x in {2.0,4.0,6.0,8.0}{
  \draw[thick, fill=white] (\x-0.14,2.4) rectangle (\x+0.14,4.6);}
\node[font=\tiny] at (2.0,2.1) {IF/ID};
\node[font=\tiny] at (4.0,2.1) {ID/EX};
\node[font=\tiny] at (6.0,2.1) {EX/MEM};
\node[font=\tiny] at (8.0,2.1) {MEM/WB};
\draw[->, very thick] (1.55,3.5)--(2.25,3.5);
\draw[->, very thick] (3.55,3.5)--(4.25,3.5);
\draw[->, very thick] (5.55,3.5)--(6.25,3.5);
\draw[->, very thick] (7.55,3.5)--(8.25,3.5);
\draw[->, dashed, accent] (9.2,2.8) .. controls (5,1.4) .. (2.9,2.8)
  node[midway, below, font=\tiny]{WB writes the regfile};
\node[align=left, text=slate, text width=11cm] at (6,0.5)
  {Four registers slice the single-cycle path into five stages.
   Each carries everything the later stages need: the instruction,
   the operands, the control bits, the pc+4. The registers are the
   whole difference --- and the whole hazard.};
\end{tikzpicture}

\caption[Five-stage datapath]{Four pipeline registers slice ch05's
path. Each register carries data \emph{and} control --- the control
bits travel with their instruction as far as they are needed.}
\label{fig:5stage06}
\end{figure}

\begin{checkpoint}[the invariant]\label{cp:invariant06}
A correct pipeline must be observationally equivalent to the
single-cycle machine: same final register file, same memory, same
order of effects, just faster. Everything in ch07 exists to
\emph{restore} that equivalence where the naive split breaks it.
\end{checkpoint}

\section{Timing diagrams}
\label{sec:timing06}
\index{timing diagram}

The convention: rows are instructions, columns are cycles, each
cell names the stage that instruction occupies. A pipeline of depth
$k$ needs $k$ cycles to fill; afterward, one completes per cycle
--- the steady state everything is measured against.

\begin{example}[fill the diagram]\label{ex:fill06}
For five instructions: cycle 1 has I1 in IF; cycle 2 has I1 in ID,
I2 in IF; \ldots; cycle 5 has I1 in WB and I5 in IF --- the first
completion; cycles 6--9 complete I2--I5. Nine cycles for five
instructions: 5 single-cycle periods would have been $5\times
5\tau$ worth of throughput 5$\times$ slower.
\end{example}

\section{Where it breaks: hazards, shown}
\label{sec:hazard06}
\index{hazard}\index{data hazard}

\begin{example}[the wrong answer, measured]\label{ex:wrong06}
\texttt{ch06\_pipe\_tb} runs two streams on the naive core. The
dependent pair \texttt{addi a0,zero,5; addi a1,a0,3}:
\begin{lstlisting}
dependent: a0=5 (want 5)  a1=3 (want 8)
  ^^ a1 read the OLD a0 -- data hazard, REAL and wrong
independent: a0..a4 = 1 2 3 4 5 (want 1..5)
ch06 SHOWN: pipeline retires ~1/cyc when independent,
lies when dependent
\end{lstlisting}
\texttt{a1=3}: the second \texttt{addi} was in ID during cycle 3,
reading the register file \emph{before} the first one's WB in
cycle 5. The machine is perfectly fast and perfectly wrong ---
and no warning prints, because hardware never warns.
\end{example}

\begin{figure}[t]
\centering
\begin{tikzpicture}[font=\scriptsize, >=latex]
% hazard timing diagram: addi a0 then addi a1,a0
\node[font=\footnotesize\bfseries] at (5.4,5.0)
  {the dependent pair, in time};
\foreach \r/\insn in {
  0/{\texttt{addi a0,zero,5}}, 1/{\texttt{addi a1,a0,3}}}{
  \node[left, font=\tiny] at (0.6,4.3-\r*0.9) {\insn};}
\foreach \i in {0,...,6}{
  \draw[thick] (1.0+\i*0.7,4.05) -- (1.0+\i*0.7,3.0)
    node[below, font=\tiny] {c\i};}
\foreach \r/\i/\st in {
  0/0/{IF}, 0/1/{ID}, 0/2/{EX}, 0/3/{MEM}, 0/4/{WB},
  1/1/{IF}, 1/2/{ID}, 1/3/{EX}, 1/4/{MEM}, 1/5/{WB}}{
  \draw[thick, fill=accentsoft] (1.0+\i*0.7,3.55-\r*0.9)
    rectangle (1.7+\i*0.7,3.95-\r*0.9);
  \node[font=\tiny] at (1.35+\i*0.7,3.75-\r*0.9) {\st};}
\draw[->, very thick, red] (4.5,3.75) -- (3.45,2.9)
  node[midway, right, font=\tiny, align=left]{a0 written\\here (WB)};
\draw[->, very thick, red] (2.4,2.9) -- (3.1,2.5);
\node[font=\tiny, red] at (2.0,2.35) {a0 read HERE (ID) --- too early};
\node[align=left, text=slate, text width=11cm] at (5.4,1.0)
  {ID for the second \texttt{addi} reads the register file two
   cycles before the first one's WB lands the new a0. The machine
   answers with the OLD value --- our tb measured a1 = 3, not 8.};
\end{tikzpicture}

\caption[Data hazard timing]{Why a1 = 3: the reader's ID lands
three cycles before the writer's WB. The register file answers
honestly --- with a stale value.}
\label{fig:hazard06}
\end{figure}

\begin{figure}[t]
\centering
\begin{tikzpicture}[font=\scriptsize, >=latex]
% CPI measured: independent vs dependent streams
\node[font=\footnotesize\bfseries] at (5,4.6)
  {measured on the naive core (real tb output)};
\draw[thick, fill=goodbg] (1,1.4) rectangle (4,3.4);
\node[font=\tiny, align=center] at (2.5,3.0) {independent stream};
\node[font=\tiny, align=center] at (2.5,2.2) {a0..a4 = 1 2 3 4 5\\
correct, ~1 retire/cycle};
\draw[thick, fill=amberbg] (5.6,1.4) rectangle (9,3.4);
\node[font=\tiny, align=center] at (7.3,3.0) {dependent stream};
\node[font=\tiny, align=center] at (7.3,2.2) {a1 = 3, not 8\\
retires 1/cycle --- wrong values};
\node[align=left, text=slate, text width=11cm] at (5,0.6)
  {The pipeline is FAST either way --- that's the trap. Throughput
   looks perfect while the dependent stream silently computes
   garbage. ch07's forwarding/stalls buy correctness back with
   real hardware.};
\end{tikzpicture}

\caption[Measured CPI outcome]{The naive core's real results:
correct and fast on independent streams, fast and wrong on
dependent ones. Throughput masks the bug --- that is what makes
hazards dangerous.}
\label{fig:cpi06}
\end{figure}

\begin{example}[name the three hazards]\label{ex:haznames06}
What we just measured is a \emph{data hazard}: an operand not yet
written. A \emph{control hazard} is the same story on the pc:
a branch resolves in EX but two wrong-path instructions are already
in IF/ID. A \emph{structural hazard} is a resource fought over:
our unified testbench memory would fight IF and MEM if it were
one port --- real machines split i-cache/d-cache partly for this.
Ch06's core takes all three on the chin; ch07 pays them off.
\end{example}

\begin{checkpoint}[ideal CPI vs honest CPI]\label{cp:cpi06}
Ideal pipelined CPI is 1. Our measured CPI is also $\approx$1 ---
on the independent stream. Predict the dependent stream's
\emph{effective} CPI if we forced correctness with stalls: each
dependent pair costs bubbles. Ch07 measures exactly this.
\end{checkpoint}

\section{Common mistakes}
\label{sec:mistakes06}
\begin{itemize}
\item \textbf{Believing throughput = correctness.} The dependent
      stream retires one per cycle too --- wrong values at full
      speed.
\item \textbf{Dropping control bits early.} A control signal that
      dies at ID/EX when MEM needs it is a silent misroute.
\item \textbf{Forgetting the fill/drain.} CPI$\approx$1 claims
      assume steady state; a 5-instruction program spends most of
      its life filling.
\item \textbf{Testing independent code only.} It passes beautifully
      and proves nothing about hazards --- our tb runs both on
      purpose.
\item \textbf{Reading the regfile on the wrong edge.} Write-first
      vs read-first ordering inside a cycle decides whether one
      hazard vanishes for free.
\end{itemize}

\section{Summary}
\label{sec:summary06}
\begin{tabularx}{\linewidth}{@{}l X@{}}
\toprule
\textbf{Idea} & \textbf{The truth} \\
\midrule
stages & IF ID EX MEM WB; 4 registers carry data+control \\
win & throughput $\times$5 ideally; latency unchanged \\
lie & retires 1/cycle even when wrong (measured: a1=3$\neq$8) \\
hazards & data, control, structural --- named, shown, deferred \\
next & forwarding, stalls, branch prediction (ch07) \\
\bottomrule
\end{tabularx}

\section*{Further reading}
\label{sec:further06}
\begin{itemize}
\item Patterson \& Hennessy, COD ch4 (pipelining, hazards):\\
      \href{https://www.elsevier.com/books/computer-organization-and-design-risc-v-edition/patterson/978-0-12-820331-6}{elsevier.com --- COD RISC-V edition}
\item Harris \& Harris, DDCA RISC-V ch7 (pipelined processor):\\
      \url{https://pages.hmc.edu/harris/ddca/ddcarv.html}
\item CS:APP ch4 (SEQ $\to$ PIPE):\\
      \url{https://csapp.cs.cmu.edu/}
\item Hennessy \& Patterson, CAQA appendix C (the classic):\\
      \href{https://www.elsevier.com/books/computer-architecture/hennessy/978-0-12-811905-1}{elsevier.com --- CAQA}
\end{itemize}

\begin{rulebox}[Watch instead]
\begin{itemize}
\item \textbf{DDCA Ch7 pipelined-processor walkthrough} --- Sarah
      Harris series, the same five-stage decomposition drawn live
      (verified 2026-10-11). \url{https://youtu.be/AoBkibslRBM}
\end{itemize}
\end{rulebox}

\section*{Notes}
\label{sec:notes06}
\begin{itemize}
\item Stage decomposition and hazard taxonomy follow COD ch4 and
      DDCA ch7; CS:APP ch4 for the SEQ$\to$PIPE pedagogy.
\item The ``a1=3'' result is a real Icarus run of
      \texttt{ch06\_pipe\_tb} --- not a stylised example. Both
      streams captured verbatim.
\end{itemize}

\section{Exercises}
\label{sec:exercises06}

\begin{exercise}[easy]\label{ex:e06a}
Latency vs throughput: which does a pipeline improve, and by how
much ideally?
\end{exercise}
\begin{solution}[hint]
Throughput, by the stage count; latency slightly worsens.
\end{solution}

\begin{exercise}[easy]\label{ex:e06b}
List what IF/ID must carry vs what MEM/WB must carry, and why the
difference.
\end{exercise}
\begin{solution}[hint]
IF/ID: insn+pc4 (everything flows forward). MEM/WB: only the
writeback payload --- control died with use.
\end{solution}

\begin{exercise}[medium]\label{ex:m06a}
Draw the timing diagram for our dependent pair and mark the exact
cells where the write and the read happen.
\end{exercise}
\begin{solution}[hint]
Write: WB cell of insn1 (c5). Read: ID cell of insn2 (c3). Two
cycles too early.
\end{solution}

\begin{exercise}[medium]\label{ex:m06b}
Why does the naive core need no flush logic yet still get branches
wrong? What would flushing cost in bubbles?
\end{exercise}
\begin{solution}[hint]
Branches resolve in EX --- two wrong-path insns already fetched;
flushing = discarding those two cells.
\end{solution}

\begin{exercise}[medium]\label{ex:m06c}
The regfile reads in ID and writes in WB. If we made writes visible
to reads in the \emph{same} cycle, which hazards disappear for free?
\end{exercise}
\begin{solution}[hint]
The WB-to-ID same-cycle case --- write-then-read ordering inside the
clock edge buys one free forwarding hop.
\end{solution}

\begin{exercise}[medium]\label{ex:m06d}
Our tb checked registers, not memory. Construct a dependent
\texttt{sw} that would store a stale value and say what signature
you'd check in dmem.
\end{exercise}
\begin{solution}[hint]
addi a0,zero,7; sw a0,0(zero) --- dmem[0] gets OLD a0; check
dmem[0]==7 fails.
\end{solution}

\begin{exercise}[medium]\label{ex:m06e}
With 5-stage delay profile $\tau$=2ns + register overhead
$\delta$=0.1ns, compute ideal speedup vs the ch05 single-cycle.
\end{exercise}
\begin{solution}[hint]
10.1ns vs 2.1ns $\to$ ~4.8$\times$, not 5 --- overhead is why real
speedup is always sublinear.
\end{solution}

\begin{exercise}[hard]\label{ex:h06a}
Derive optimal stage count minimax-style: total speedup =
$T/((T/k)+d)$. Find $k$ that maximises it for $T$=10, $d$=0.1.
\end{exercise}
\begin{solution}[hint]
Differentiate or tabulate: $k=\sqrt{T/d}\approx10$ stages; past
that, register overhead eats the gain.
\end{solution}

\begin{exercise}[hard]\label{ex:h06b}
Three instructions: \texttt{lw a0,0(x); add a1,a0,a2; sw a1,4(x)}.
On the naive core trace every wrong register read and say how many
independent nops would fix each.
\end{exercise}
\begin{solution}[hint]
a0's load lands WB at c5; add reads at c3 $\to$ 2 nops. Then a1 is
an ALU result, WB c6, read c4 $\to$ 2 more nops.
\end{solution}

% chapter end
